



















\subsection{行列式的性质}
\index{determinant|(}
我们想要一个公式，
用以判断一个 $\nbyn{n}$ 矩阵是否非奇异的。
我们不打算一开始就陈述这样一个公式，
而是先针对每个~$n$，
考虑这样的公式所计算的函数。
我们将通过一列性质来定义这个函数。
随后我们将证明具有这些性质的函数存在且唯一，
并说明如何计算它。
（由于我们最终会证明这一点，所以从一开始
我们就直接说“\( \det(T) \)”而不说
“若存在唯一的行列式函数，则 \( \det(T) \)”。）

\begin{definition}  \label{def:Det}
%<*df:Det>
一个 \( \nbyn{n} \) \definend{行列式\/}\index{determinant!definition}%
\index{matrix!determinant}是一个函数
\( \map{\det}{\matspace_{\nbyn{n}}}{\Re} \)，满足下列条件：
\begin{enumerate}
  \item
    $
       \det (\vec{\rho}_1,\dots,k\cdot\vec{\rho}_i
                                     + \vec{\rho}_j,\dots,\vec{\rho}_n)
       =\det (\vec{\rho}_1,\dots,\vec{\rho}_j,\dots,\vec{\rho}_n)
    $ 当 \( i\ne j \) 时
  \item
    $
       \det (\vec{\rho}_1,\ldots,\vec{\rho}_j,
               \dots,\vec{\rho}_i,\dots,\vec{\rho}_n)
       = -\det (\vec{\rho}_1,\dots,\vec{\rho}_i,\dots,\vec{\rho}_j,
                \dots,\vec{\rho}_n)
    $
    当 \( i\ne j\) 时
  \item
    $
       \det (\vec{\rho}_1,\dots,k\vec{\rho}_i,\dots,\vec{\rho}_n)
       = k\cdot \det (\vec{\rho}_1,\dots,\vec{\rho}_i,\dots,\vec{\rho}_n)
    $
    对任意标量~$k$
    % for \( k\ne 0\)
  \item
     $
       \det(I)=1
     $
     其中 \( I \) 是单位矩阵
\end{enumerate}
（各个 $\vec{\rho}\,$ 是矩阵的行）。
我们常把 \( \det (T) \) 写成 \( \deter{T} \)。
%</df:Det>
\end{definition}

\begin{remark}  \label{rem:SwapRowsRedun}
%<*re:SwapRowsRedun>
条件~(2) 是多余的，因为
 \begin{equation*}
   T\grstep{\rho_i+\rho_j}
    \repeatedgrstep{-\rho_j+\rho_i}
    \repeatedgrstep{\rho_i+\rho_j}
    \repeatedgrstep{-\rho_i}
    \hat{T}
\end{equation*}
对换第 \( i \) 行与第~\( j \) 行。
我们列出它是为了
与前面各章中高斯消元法的表述保持一致。
%</re:SwapRowsRedun>
\end{remark}

\begin{remark}
条件~(3) 没有 $k\neq 0$ 的限制，虽然
高斯消元法中把一行乘以~$k$ 的操作
确有此限制。
% \nearbylemma{le:IdenRowsDetZero} below
下面的结果表明，在这里我们不需要该限制。
\end{remark}

% % \begin{remark}
% The previous subsection's plan
% asks for the determinant of an echelon form matrix to be the product
% down the diagonal.
% The next result shows that in the presence of the other
% three, condition~(4) gives that.
% % \end{remark}

\begin{lemma}   \label{le:IdenRowsDetZero}
%<*lm:IdenRowsDetZero>
有两行相同的矩阵，其行列式为零。
有一个零行的矩阵，其行列式为零。
矩阵是非奇异的当且仅当其行列式非零。
阶梯形矩阵的行列式等于沿其对角线的乘积。
%</lm:IdenRowsDetZero>
\end{lemma}

\begin{proof}
%<*pf:IdenRowsDetZero0>
为验证第一句，对换这两个相同的行。
行列式的符号改变了，但矩阵是同一个，
因而其行列式也相同。
于是行列式为零。
%</pf:IdenRowsDetZero0>

%<*pf:IdenRowsDetZero1>
对于第二句，
将零行乘以二。
这使行列式加倍，但该行
保持不变，因而
行列式也保持不变。
所以行列式必为零。
%</pf:IdenRowsDetZero1>

%<*pf:IdenRowsDetZero2>
对于第三句，
作高斯–若尔当消元
$T \rightarrow\cdots\rightarrow\hat{T}$。
由前三个性质，
$T$ 的行列式为零当且仅当
$\hat{T}$ 的行列式为零
（尽管二者可能在符号或大小上不同）。
非奇异矩阵 $T$ 经高斯–若尔当消元化为单位矩阵，
因而其行列式非零。
奇异的 $T$ 则化为带有一个零行的 $\hat{T}$；
由本引理的第二句，其行列式为零。
%</pf:IdenRowsDetZero2>

%<*pf:IdenRowsDetZero3>
第四句有两种情形。
若阶梯形矩阵是奇异的，则
它有一个零行。
于是其对角线上有一个零，并且
沿其对角线的乘积为零。
由本结论的第三句，行列式为零，
因此该矩阵的行列式等于
沿其对角线的乘积。
%</pf:IdenRowsDetZero3>

%<*pf:IdenRowsDetZero4>
若阶梯形矩阵是非奇异的，则它的对角元
都不为零。
这意味着我们可以除以这些对角元，
再利用条件~(3) 使对角线上都得到 $1$。
\begin{equation*}
  \begin{vmat}
    t_{1,1}  &t_{1,2}  &     &t_{1,n}  \\
    0        &t_{2,2}  &     &t_{2,n}  \\
             &         &\ddots         \\
    0        &         &     &t_{n,n}
  \end{vmat}
  =
  t_{1,1}\cdot t_{2,2}\cdots t_{n,n}\cdot
  \begin{vmat}
    1        &t_{1,2}/t_{1,1}  &     &t_{1,n}/t_{1,1}  \\
    0        &1                &     &t_{2,n}/t_{2,2}  \\
             &                 &\ddots         \\
    0        &                 &     &1
  \end{vmat}
\end{equation*}
然后，高斯–若尔当消元的若尔当部分使之成为单位矩阵。
\begin{equation*}
  =
  t_{1,1}\cdot t_{2,2}\cdots t_{n,n}\cdot
  \begin{vmat}
    1        &0                &     &0                \\
    0        &1                &     &0                \\
             &                 &\ddots         \\
    0        &                 &     &1
  \end{vmat}
  =
  t_{1,1}\cdot t_{2,2}\cdots t_{n,n}\cdot 1
\end{equation*}
所以在这种情形，行列式
同样是沿对角线的乘积。
%</pf:IdenRowsDetZero4>
\end{proof}

这给了我们一种办法，用来计算行列式函数
在一个矩阵上的取值：
作高斯消元，记录行对换所引起的任何变号
以及我们提出的任何标量倍数，
最后沿阶梯形结果的
对角线相乘。
这一算法与高斯消元法一样快，因此
对我们将会看到的
所有矩阵都是实用的。

\begin{example}
用高斯消元法
计算 \( \nbyn{2} \) 行列式
\begin{equation*}
   \begin{vmat}[r]
      2  &4  \\
      -1 &3
   \end{vmat}
   =
   \begin{vmat}[r]
      2  &4  \\
      0  &5
   \end{vmat}
   =10
\end{equation*}
省不了多少时间，
因为 $\nbyn{2}$ 行列式公式很容易。
然而，一个 \( \nbyn{3} \) 行列式
用高斯消元法来计算往往比用它的公式更容易。
\begin{equation*}
   \begin{vmat}[r]
     2  &2  &6  \\
     4  &4  &3  \\
     0  &-3 &5
   \end{vmat}
   =
   \begin{vmat}[r]
     2  &2  &6  \\
     0  &0  &-9 \\
     0  &-3 &5
   \end{vmat}
   =
   -\begin{vmat}[r]
     2  &2  &6  \\
     0  &-3 &5  \\
     0  &0  &-9
   \end{vmat}
   =-54
\end{equation*}
\end{example}

\begin{example}
比 $\nbyn{3}$ 更大的行列式
用高斯消元法的过程很快就完成了。
\begin{equation*}
   \begin{vmat}[r]
      1  &0  &1  &3  \\
      0  &1  &1  &4  \\
      0  &0  &0  &5  \\
      0  &1  &0  &1
   \end{vmat}
   =
   \begin{vmat}[r]
      1  &0  &1  &3  \\
      0  &1  &1  &4  \\
      0  &0  &0  &5  \\
      0  &0  &-1 &-3
   \end{vmat}
   =
   -\begin{vmat}[r]
      1  &0  &1  &3  \\
      0  &1  &1  &4  \\
      0  &0  &-1 &-3 \\
      0  &0  &0  &5
   \end{vmat}
   =-(-5)=5
\end{equation*}
\end{example}

上面的例题引出了一个重要的观点。
本章的引言给出了 $\nbyn{2}$
和 $\nbyn{3}$ 行列式的公式，因此我们知道它们存在，
但对于 $\nbyn{4}$ 或更大的矩阵上的行列式函数却没有给出公式。
作为替代，\nearbydefinition{def:Det} 给出了一个行列式函数
应当具有的性质，
并由此引出用高斯消元法来计算行列式。

然而，对任意一个矩阵，我们都可以用高斯消元法以
多种方式将其化为阶梯形。
例如，给定一个消元过程，我们可以插入一个第一步
把最上面一行乘以~$2$，然后插入一个第二步
把它乘以~$1/2$，从而改变这个过程。
因此我们必须担心，两种不同的高斯消元法消元
可能导致行列式的两个不同的计算值。

也就是说，我们必须验证 \nearbydefinition{def:Det} 给出的是
一个良定义的函数。
接下来的两小节将完成这件事，证明存在
满足该定义的
一个良定义的函数。

但首先我们证明，若存在这样的函数，
则至多只有一个。
上面的例题说明了这个想法：~我们
按照定义的性质得到了~$5$。
因此，虽然我们尚未证明 $\det_{\nbyn{4}}$ 存在，
即存在具有性质 (1)\,--\,(4) 的函数，
但若满足这些性质的函数确实存在，
则我们知道它在上述矩阵上给出的值。

\begin{lemma} \label{lm:DetFcnIsUnique}
%<*lm:DetFcnIsUnique>
对每个 $n$，
若存在 $\nbyn{n}$ 行列式函数，则它是唯一的。
%</lm:DetFcnIsUnique>
\end{lemma}

\begin{proof}
%<*pf:DetFcnIsUnique>
假设有两个函数
$\map{\det_1,\det_2}{\matspace_{\nbyn{n}}}{\Re}$
满足
\nearbydefinition{def:Det} 的性质以及
其推论 \nearbylemma{le:IdenRowsDetZero}。
给定一个方阵~$M$，固定一种实施高斯消元法
将该矩阵化为阶梯形的方式（有多种方式并不要紧，
只需固定其中一种）。
通过像上面的例题中那样
使用这一固定的消元过程\Dash
记录行倍乘的因子以及符号在行对换时如何变号，
然后沿阶梯形结果的对角线相乘\Dash
我们可以算出这两个函数在~$M$ 上必定返回的值，
并且它们必定返回相同的值。
由于它们在每个输入上都给出相同的输出，所以它们是同一个函数。
%</pf:DetFcnIsUnique>
\end{proof}

“若存在 $\nbyn{n}$ 行列式函数”
这一说法强调，
虽然我们可以
用高斯消元法计算出行列式函数
所可能返回的唯一值，
但我们尚未证明这样的函数对一切 $n$ 都存在。
本节其余部分将完成这件事。



